\lysh{38807}{4}

\begin{alist}
\item Το πιεσόμετρο έχει απόκλιση $-4\text{mmHg}$ (δίνει τιμές $4\text{mmHg}$ μικρότερες από τις πραγματικές).
Η πραγματική τιμή για τη διαστολική πίεση είναι: $p = p^\prime + 4$ και η πραγματική τιμή για τη συστολική πίεση είναι: $P = P^\prime + 4$.

\item Θεωρούμε ότι έχουμε υπόταση, αν είναι $P < 90$ και $p < 60$. Αντικαθιστούμε τις σχέσεις $P = P^\prime + 4$ και $p = p^\prime + 4$:
$P^\prime + 4 < 90$ και $p^\prime + 4 < 60$,
άρα $P^\prime < 86$ και $p^\prime < 56$, οπότε έχουμε υπόταση, όταν οι μετρήσεις με το πιεσόμετρο δείξουν τιμές κάτω από αυτές.

Θεωρούμε ότι έχουμε υπέρταση, αν $P > 140$ και $p > 90$. Αντικαθιστούμε τις σχέσεις:
$P^\prime + 4 > 140$ και $p^\prime + 4 > 90$,
άρα $P^\prime > 136$ και $p^\prime > 86$, οπότε έχουμε υπέρταση, όταν το πιεσόμετρο δώσει μέτρηση πάνω από τις τιμές αυτές.

\item
\begin{rlist}
\item Η σχέση θα είναι: $40 \le \text{συστολική} - \text{διαστολική} \le 60$, δηλαδή
$40 \le P^\prime - p^\prime \le 60$.

\item Θα λύσουμε την ανίσωση για $p^\prime = 80$:
$40 \le P^\prime - 80 \le 60 \Leftrightarrow 40 + 80 \le P^\prime \le 60 + 80$, άρα $120 \le P^\prime \le 140$.
\end{rlist}
\end{alist}
